\section{Conclusions}
\label{sec:conclusions}

RI-selected heads participate in causal paths supporting contextual mother-of
retrieval, but RI poorly identifies the heads that matter. It misses major
contributors and selects many heads with weak effects. The mechanistic evidence
instead supports interacting writers and readers: some RI heads receive fact-site
information and influence later answer-position computation. Their roles include
both promoting and opposing the answer; copying alone does not explain them.
This establishes participation in relation use, not that individual heads encode
the relation itself (\S\S\ref{sec:results-why}--\ref{sec:results-routes}).

Weak standalone effects do not imply irrelevance. L9H16 and L1H27 improve
fact--query discrimination within the partial retained background, with effects
reproduced on held-out families. Weaker RI heads also matter collectively after
excluding the eight strong members. These findings establish context-dependent
contributions, but not a general hidden-backup mechanism: the individual effects
depend strongly on replacement baseline, and no new route assignment passes the
discovery rule. The selected 50-head background also fails to preserve the task,
leaving a sufficient circuit unresolved.

Across training, early routing changes do not identify the mature high-RI heads
(Appendix~\ref{app:dev}); developmental association likewise cannot substitute
for causal evidence.
